% !TeX program = xelatex
% !TeX encoding = UTF-8
% 编译方法：xelatex 合并稿.tex（建议连续编译两次以更新交叉引用）
\documentclass[UTF8,a4paper,zihao=-4,fontset=none]{ctexart}

% 字体采用三级自动选择机制，使模板能够在不同电脑上直接编译。
% 先检查 Windows 标准字体目录；再检查 macOS 版 Microsoft Word 的字体；均不可用时回退到 TeX Live 自带的 Fandol 字体。
% 无论选择哪一组字体，主字体的 BoldFont 都显式绑定到对应黑体，确保 \textbf 下的中文字体正确。
\newcommand{\UseWindowsMicrosoftFonts}{
  \setCJKmainfont{simsun.ttc}[Path={C:/Windows/Fonts/},BoldFont=simhei.ttf]
  \setCJKsansfont{simhei.ttf}[Path={C:/Windows/Fonts/},BoldFont=simhei.ttf]
  \setCJKmonofont{simsun.ttc}[Path={C:/Windows/Fonts/}]
  \setCJKfamilyfont{paperhei}{simhei.ttf}[Path={C:/Windows/Fonts/},BoldFont=simhei.ttf]
  \PackageInfo{paper-template}{Using Windows SimSun and SimHei}
}

\newcommand{\UseOfficeMicrosoftFonts}{
  \setCJKmainfont{Simsun.ttc}[
    Path={/Applications/Microsoft Word.app/Contents/Resources/DFonts/},
    BoldFont=SimHei.ttf
  ]
  \setCJKsansfont{SimHei.ttf}[
    Path={/Applications/Microsoft Word.app/Contents/Resources/DFonts/},
    BoldFont=SimHei.ttf
  ]
  \setCJKmonofont{Simsun.ttc}[
    Path={/Applications/Microsoft Word.app/Contents/Resources/DFonts/}
  ]
  \setCJKfamilyfont{paperhei}{SimHei.ttf}[
    Path={/Applications/Microsoft Word.app/Contents/Resources/DFonts/},
    BoldFont=SimHei.ttf
  ]
  \PackageInfo{paper-template}{Using Microsoft Office SimSun and SimHei}
}

\newcommand{\UsePortableFandolFonts}{
  \setCJKmainfont{FandolSong-Regular.otf}[cmap=UniGB-UTF16-H,BoldFont=FandolHei-Regular.otf]
  \setCJKsansfont{FandolHei-Regular.otf}[cmap=UniGB-UTF16-H,BoldFont=FandolHei-Regular.otf]
  \setCJKmonofont{FandolFang-Regular.otf}[cmap=UniGB-UTF16-H]
  \setCJKfamilyfont{paperhei}{FandolHei-Regular.otf}[cmap=UniGB-UTF16-H,BoldFont=FandolHei-Regular.otf]
  \PackageWarningNoLine{paper-template}{SimSun and SimHei not found; using portable Fandol fonts}
}

\newcommand{\UseOfficeOrPortableFonts}{
  \IfFileExists{/Applications/Microsoft Word.app/Contents/Resources/DFonts/Simsun.ttc}{
    \IfFileExists{/Applications/Microsoft Word.app/Contents/Resources/DFonts/SimHei.ttf}{
      \UseOfficeMicrosoftFonts
    }{
      \UsePortableFandolFonts
    }
  }{
    \UsePortableFandolFonts
  }
}

% 将下面的 0 改为 1，可强制使用完全随 TeX Live 提供的 Fandol 字体进行兼容性测试。
\providecommand{\ForcePortableFonts}{0}
\ifnum\ForcePortableFonts=1
  \UsePortableFandolFonts
\else
  \IfFileExists{C:/Windows/Fonts/simsun.ttc}{
    \IfFileExists{C:/Windows/Fonts/simhei.ttf}{
      \UseWindowsMicrosoftFonts
    }{
      \UseOfficeOrPortableFonts
    }
  }{
    \UseOfficeOrPortableFonts
  }
\fi
\newcommand{\heifont}{\CJKfamily{paperhei}}

% 页面、行距与段落。
\usepackage[a4paper,top=2.5cm,bottom=2.5cm,left=2.5cm,right=2.5cm,footskip=1.2cm]{geometry}
\usepackage{setspace}
\singlespacing
\setlength{\parindent}{2em}
\setlength{\parskip}{0pt}
\setlength{\emergencystretch}{3em}
\usepackage{indentfirst}

% 数学公式、表格与插图。
\usepackage{amsmath,amssymb,mathtools,bm}
\usepackage{graphicx}
\graphicspath{{./}}
\usepackage{booktabs,tabularx,array,multirow}
\usepackage[table]{xcolor}
\usepackage{float}
\usepackage{caption}
\usepackage{subcaption}
\usepackage{enumitem}

% 标题、页码与交叉引用。
\usepackage{titlesec}
\usepackage{fancyhdr}
\usepackage[hidelinks]{hyperref}
\usepackage{cleveref}

% 一级标题使用中文序号并居中，下级标题使用阿拉伯数字层级编号并左对齐。
\renewcommand{\thesection}{\arabic{section}}
\renewcommand{\thesubsection}{\arabic{section}.\arabic{subsection}}
\renewcommand{\thesubsubsection}{\arabic{section}.\arabic{subsection}.\arabic{subsubsection}}
\setcounter{secnumdepth}{3}
\titleformat{\section}[block]
  {\centering\heifont\bfseries\zihao{4}}
  {\chinese{section}、}{0pt}{}
\titleformat{\subsection}[hang]
  {\raggedright\heifont\bfseries\zihao{-4}}
  {\thesubsection}{0.8em}{}
\titleformat{\subsubsection}[hang]
  {\raggedright\heifont\bfseries\zihao{-4}}
  {\thesubsubsection}{0.8em}{}
\titlespacing*{\section}{0pt}{2.8ex plus 0.5ex minus 0.2ex}{1.7ex}
\titlespacing*{\subsection}{0pt}{1.6ex plus 0.3ex minus 0.1ex}{0.8ex}
\titlespacing*{\subsubsection}{0pt}{1.4ex plus 0.2ex minus 0.1ex}{0.7ex}

% 页码从摘要页开始，位于页脚中央。
\pagestyle{fancy}
\fancyhf{}
\fancyfoot[C]{\zihao{5}\thepage}
\renewcommand{\headrulewidth}{0pt}
\renewcommand{\footrulewidth}{0pt}

% 图表及列表的常用设置。
\captionsetup{font=small,labelfont=bf,labelsep=quad}
\setlist[itemize]{itemsep=0.35\baselineskip,topsep=0.8ex,parsep=0pt,partopsep=0pt}
\newcolumntype{C}[1]{>{\centering\arraybackslash}p{#1}}
\newcolumntype{L}[1]{>{\raggedright\arraybackslash}p{#1}}
\newcolumntype{Y}{>{\raggedright\arraybackslash}X}
\newcolumntype{Z}{>{\centering\arraybackslash}X}

% 请优先修改下面三个字段。
\newcommand{\papertitle}{无线电干扰源的快速自动定位与清除}
\newcommand{\paperabstract}{}
\newcommand{\paperkeywords}{}

\hypersetup{
  pdftitle={\papertitle},
  pdfauthor={}
}

\begin{document}
\pagenumbering{arabic}
\setcounter{page}{1}
\thispagestyle{fancy}

% 第一页：标题、摘要与关键词。
\begin{center}
  \vspace*{0.3cm}
  {\heifont\bfseries\zihao{3}\papertitle\par}
  \vspace{1.5\baselineskip}
  {\heifont\bfseries\zihao{4}摘要\par}
\end{center}

\paperabstract\par
\vspace{0.5\baselineskip}
\noindent{\heifont\bfseries 关键词：}\paperkeywords

\newpage

\section{问题背景与重述}

\subsection{问题背景}

无线电干扰源的定位与清除是无线电频谱管理的重要任务。实际排查通常经历初步监测、区域性排查和精准定位三个阶段，其中前两个阶段确定干扰源的分布区域，而具体位置仍需进一步查明。受环境遮挡、干扰源发射功率等因素影响，精准定位阶段需要技术人员携带便携式测向设备进入目标区域逐步排查，并在确认位置后完成清除。为了尽快消除干扰，需要尽可能缩短这一过程的耗时。

针对危险环境下的排查需求，某公司拟研发搭载于机器狗的干扰源自动定位清除系统。该系统通过无线电测向获取干扰源的方位信息，并在接近目标后利用光学探测设备实现精确定位和清除。然而，测向结果存在误差，信号接收受到距离和发射方向的限制，移动与各项设备操作也需要耗时。因此，如何在有限观测信息下完成干扰源的搜索与定位，并在确保全部清除的前提下提高作业效率，是本题研究的核心问题。

\subsection{问题重述}

题目已将干扰源的分布范围确定为半径 \(1800\,\mathrm{m}\) 的圆形区域，以圆心为坐标原点，正东方向和正北方向分别为 \(x\) 轴与 \(y\) 轴正向。各干扰源占用互不相同且保持不变的频道，频道编号为 \(1\) 至 \(20\)，不同干扰源的信号互不影响。干扰源分为全向与定向两类，其有效覆盖角度分别为 \(360^\circ\) 和定向方向两侧各 \(90^\circ\)，有效接收半径在 \(1000\,\mathrm{m}\) 至 \(1500\,\mathrm{m}\) 之间。

测向机每次只能检测一个频道，测得的示向度与真实方位角之间的误差在 \([-1^\circ,1^\circ]\) 内，同一地点在一段时间内重复检测不会改变该误差。机器狗以 \(5\,\mathrm{m/s}\) 的速度沿直线移动，且必须停止移动后才能有效检测。频道切换和单次检测分别耗时 \(1\,\mathrm{s}\) 和 \(5\,\mathrm{s}\)；当机器狗与干扰源的距离不超过 \(20\,\mathrm{m}\) 时，可依次进行光学精确定位和清除，两项操作分别耗时 \(3\,\mathrm{s}\) 和 \(2\,\mathrm{s}\)。依据上述条件，需要完成以下四项任务。

\textbf{问题一：}已知若干检测点的坐标及其对同一干扰源测得的示向度，给出采用交会定位法计算多边形定位区域直径的算法，并判断以该直径为直径的圆能否覆盖整个定位区域。

\textbf{问题二：}已知某全向干扰源在一个检测点处的示向度，提出第二个检测点的选择策略，使进一步检测能够取得较好的定位效果，并给出第二个检测点的候选区域。

\textbf{问题三：}目标区域内有 \(10\) 至 \(16\) 个全向干扰源，具体数量未知。机器狗从原点出发，测向机初始频道为 \(1\)，需要自主完成搜索、定位和清除，并在确保所有干扰源被清除的前提下尽可能缩短任务时间。制定相应策略与算法，通过模拟器演练，统计被清除干扰源个数的比例及平均定位清除时间，以检验任务完成情况与作业效率。

\textbf{问题四：}目标区域内同时存在全向和定向干扰源，总数仍为 \(10\) 至 \(16\) 个，但具体数量、定向干扰源的数量及其定向方向均未知。其余条件与问题三相同，在此情形下提出检测与清除策略及算法，并通过模拟器演练检验其有效性。

对问题三和问题四，在充分演练后分别完成三次正式测试，并按题目规定记录测试结果，并将对应日志文件纳入支撑材料。

\section{模型假设}

为使各问的建模前提清楚，本文区分测角误差的有界约束与期望模型另行引入的概率假设，并说明数值近似的适用范围。

\begin{itemize}[label=\textbullet,leftmargin=2em]
  \item \textbf{假设一（测角误差有界）。} 问题一及问题二的最坏情形模型只使用误差上界 \(\alpha=\pi/180\)，不要求误差服从特定分布。题面给出了 \([-1^\circ,1^\circ]\) 的误差范围，并未给出其概率分布；问题二的期望模型所需的均匀分布与独立性假设，仅在相应小节内额外引入。
  \item \textbf{假设二（圆盘约束的保守近似）。} 问题一的数值实现若需将圆盘约束转为多边形约束，采用外接正多边形，以保证不因离散化排除真实位置。该近似可与半平面约束统一处理，但计算结果可能偏保守；问题二的理论区域仍按原始圆盘与距离条件定义。
  \item \textbf{假设三（同点误差不因重复而减小）。} 在同一位置、题面所述误差保持不变的时段内，重复检测同一频道不产生新的独立测角约束，不通过重复次数缩小误差界。否则会高估重复测量的信息收益；问题二对不同检测点引入的概率假设不改变这一规定。
\end{itemize}

\section{符号说明}

\begin{table}[H]
  \centering
  \begingroup
  \renewcommand{\arraystretch}{1.45}
  \begin{tabularx}{\textwidth}{C{0.18\textwidth}ZC{0.18\textwidth}}
    \toprule[1.2pt]
    {\heifont\bfseries 符号} & {\heifont\bfseries 定义} & {\heifont\bfseries 单位} \\
    \midrule[0.5pt]
    \(\Omega\) & 以原点为圆心、半径为 \(1800\,\mathrm{m}\) 的目标圆域 & -- \\
    \(g\) & 干扰源的位置，固定但未知 & \(\mathrm{m}\) \\
    \(s_i\) & 第 \(i\) 个检测点的位置 & \(\mathrm{m}\) \\
    \(q\) & 待选择的检测点或行动位置 & \(\mathrm{m}\) \\
    \(R\) & 干扰源的固定有效接收半径，取值为 \(1000\)--\(1500\,\mathrm{m}\) & \(\mathrm{m}\) \\
    \(\theta_i\) & 在检测点 \(s_i\) 获得的正常示向度读数 & \(\mathrm{rad}\) \\
    \(\alpha\) & 测角误差上界，\(\alpha=\pi/180\) & \(\mathrm{rad}\) \\
    \(r_{\mathrm{strong}}\) & 强信号距离阈值，取值为 \(5\,\mathrm{m}\) & \(\mathrm{m}\) \\
    \(r_{\mathrm{opt}}\) & 光学精确定位距离，取值为 \(20\,\mathrm{m}\) & \(\mathrm{m}\) \\
    \(\mathcal W_i\) & 第 \(i\) 次正常测向对应的正向角域，半角为 \(\alpha\) & -- \\
    \(P\) & 纯测角交会区域，即各角域 \(\mathcal W_i\) 的交集 & -- \\
    \(K\) & 综合观测信息与物理条件得到的源位置可行区域 & -- \\
    \(\operatorname{diam}(K)\) & 区域 \(K\) 的直径 & \(\mathrm{m}\) \\
    \(c(K)\) & 区域 \(K\) 的最小包围圆圆心 & \(\mathrm{m}\) \\
    \(\rho(K)\) & 区域 \(K\) 的最小包围圆半径 & \(\mathrm{m}\) \\
    \(d_{\max}(q;K)\) & 位置 \(q\) 到区域 \(K\) 的最远距离 & \(\mathrm{m}\) \\
    \bottomrule[1.2pt]
  \end{tabularx}
  \endgroup
\end{table}

上述位置变量均为二维坐标向量；公式推导中的角度统一采用弧度，观测数据与结果展示可采用度并注明单位。多源情形下增加源编号下标，其余局部符号在各问建模时另行定义。

\section{问题分析}
\subsection{问题一分析}

问题一研究使用交会定位法排查干扰源时定位区域的几何性质。根据题意，可将各次观测对应的方向范围表示为多个半平面约束，并将定位区域表示为其交集，进而利用凸多边形的性质，将区域直径的计算转化为顶点间最大距离的求解。对于同直径圆的覆盖问题，可从最远顶点对对圆心位置的限制入手，结合最小包围圆半径与区域直径的关系建立判据，并构造符合测向条件的反例检验其普遍性。

\subsection{问题二分析}

问题二将关注点从已有检测结果转向下一次检测的位置选择。此时仅掌握一个检测点的示向度，干扰源的距离和准确位置仍未知，第二个检测点必须依据当前有限的信息选取。考虑到干扰源定位精度的重要指标是定位区域最小包围圆的半径，本文将这一变量的最小化作为寻找候选点模型的优化目标。
\subsection{问题三分析}


\subsection{问题四分析}


\section{问题一的模型建立与求解}

\subsection{单次观测的半平面表示}

设共有 \(n\) 次正常示向度观测，沿用符号表中的误差上界 \(\alpha=\pi/180\)。本问用 \(x\in\mathbb R^2\) 表示候选源位置，记 \(u(\varphi)=(\cos\varphi,\sin\varphi)\) 为方向角 \(\varphi\) 对应的单位向量，并定义二维叉积 \([a,b]=a_xb_y-a_yb_x\)。一次示向度观测对应正向角域
\begin{equation}
  \mathcal W_i=\{x:[u(\theta_i-\alpha),\,x-s_i]\ge0,\quad [u(\theta_i+\alpha),\,x-s_i]\le0\}.
  \label{eq:wedge}
\end{equation}
其中第一个约束要求目标位于下边界射线的左侧，第二个约束要求目标位于上边界射线的右侧。由于 \(0<\alpha<\pi/2\)，两个半平面之交即为前向角域，排除了检测点背后的扇区。

记 \(u_i^-=(\cos(\theta_i-\alpha),\sin(\theta_i-\alpha))\)、\(u_i^+=(\cos(\theta_i+\alpha),\sin(\theta_i+\alpha))\)。在一次观测对应的两条边界上，
\begin{equation}
  \begin{bmatrix}u_{iy}^- & -u_{ix}^-\\[2pt] -u_{iy}^+ & u_{ix}^+\end{bmatrix}x
  \;\le\;
  \begin{bmatrix}u_{iy}^- & -u_{ix}^-\\[2pt] -u_{iy}^+ & u_{ix}^+\end{bmatrix}s_i .
  \label{eq:halfplane}
\end{equation}
将全部观测的系数按行组成矩阵 \(A\)，右端项组成向量 \(b\)，叠加得到
\begin{equation}
  P=\bigcap_{i=1}^{n}\mathcal W_i=\{x:Ax\le b\}.
  \label{eq:P}
\end{equation}
该表示在 \(90^\circ\)、\(270^\circ\) 以及 \(0^\circ/360^\circ\) 附近均可用，这是选择叉积形式而非斜率形式的主要原因 \cite{Tokekar2013}。若真实目标落在每个 \(\mathcal W_i\) 内，则真实位置 \(g\in P\)。\(P\) 是闭凸多面集，但不保证有界。

本问主模型始终取 \(\alpha=\pi/180\)；工程敏感性检验若将误差半宽放宽为 \(1.005^\circ\)，相应取 \(\alpha=1.005\pi/180\)，并与主模型结果分别记录。

\subsection{区域状态判定与顶点算法}

算法的处理步骤如下。

\begin{enumerate}[label=(\arabic*),leftmargin=2.2em]
  \item 归一化各半平面法向量，求解可行性线性规划。若不可行则返回 empty，不得继续输出伪定位点。
  \item 将候选位置写作 \(x=(x_1,x_2)\)，在 \(P\) 上分别最大化 \(x_1\)、\(-x_1\)、\(x_2\)、\(-x_2\)。任一目标无界则 \(P\) 无界，直径记为 \(+\infty\)；四个方向均有有限上界则 \(P\) 有界。
  \item 取一个可行点作坐标平移，减轻大坐标下的数值相消。
  \item 枚举不平行的边界直线对，求其交点，仅保留满足全部半平面约束的交点，并按容差去重。
  \item 对二维区域按极角整理顶点顺序；对点、线段分别保留退化表示。
  \item 输出原约束的最大残差，并计算直径与最小包围圆，给出状态诊断。
\end{enumerate}

若共有 \(m=2n\) 条约束，交点枚举与可行性检验的复杂度为 \(O(m^3)\)，顶点存储为 \(O(h)\)（\(h\) 为去重后的顶点数）。

若出现空交集，意味着出现了观测不一致的情况，可能来自频道误关联、角度单位错误、舍入误差或数值问题，需保留并输出诊断。

\subsection{定位区域直径}

若 \(P\) 非空有界，将其 \(h\) 个顶点记为 \(v_1,\dots,v_h\)。沿用符号表中的直径记法，有
\begin{equation}
  \operatorname{diam}(P)=\max_{x,y\in P}\lVert x-y\rVert=\max_{i,j}\lVert v_i-v_j\rVert .
  \label{eq:diameter}
\end{equation}

\textbf{证明。} 任一 \(x\in P\) 可写为顶点凸组合 \(x=\sum_i\mu_iv_i\)（\(\mu_i\ge0\)，\(\sum_i\mu_i=1\)），同理 \(y=\sum_j\nu_jv_j\)。由三角不等式，
\begin{equation}
  \lVert x-y\rVert
  =\Bigl\lVert\sum_{i,j}\mu_i\nu_j(v_i-v_j)\Bigr\rVert
  \le\sum_{i,j}\mu_i\nu_j\lVert v_i-v_j\rVert
  \le\max_{i,j}\lVert v_i-v_j\rVert .
\end{equation}
又因顶点本身属于 \(P\)，上界可以取到，故式 \eqref{eq:diameter} 成立。\(\square\) 
这一“由可行多面体顶点确定最坏位置误差”的结构与文献 \cite{Gholami2015} 的结论一致。算法上枚举 \(h(h-1)/2\) 个点对，复杂度 \(O(h^2)\)。

\subsection{定位区域的覆盖}

沿用符号表中的最小包围圆半径记法，对区域 \(P\) 有
\begin{equation}
  \rho(P)=\min_c\max_{x\in P}\lVert x-c\rVert
  =\min_c\max_i\lVert v_i-c\rVert .
  \label{eq:mincircle}
\end{equation}
第二个等号与式 \eqref{eq:diameter} 同理，是凸组合与范数凸性的结果。该最小化问题的最优圆心记为 \(c(P)\)。由于覆盖圆必须包含最远点对，有 \(\rho(P)\ge\operatorname{diam}(P)/2\)。

\textbf{判定准则。} 存在直径恰为 \(\operatorname{diam}(P)\) 的覆盖圆，当且仅当 \(\rho(P)=\operatorname{diam}(P)/2\)。

\textbf{证明。} 取任意最远点对 \(a,b\)，\(\lVert a-b\rVert=\operatorname{diam}(P)\)。若存在半径 \(\operatorname{diam}(P)/2\) 的圆（圆心 \(c\)）同时包含 \(a\) 与 \(b\)，则
\begin{equation}
  \operatorname{diam}(P)=\lVert a-b\rVert\le\lVert a-c\rVert+\lVert b-c\rVert\le\operatorname{diam}(P) .
\end{equation}
上式所有不等号必须同时取等，这要求 \(c=(a+b)/2\) 且 \(a,b\) 都在圆周上。因此只需检查所有顶点到中点 \((a+b)/2\) 的距离是否均不超过 \(\operatorname{diam}(P)/2\)。\(\square\)

下面构造一个确实能由三次带 \(\pm1^\circ\) 误差的示向度观测产生的定位区域。

取边长 \(L=38\) 米的等边三角形 \(T\)，顶点为
\begin{equation}
  v_1=(0,0),\quad v_2=(38,0),\quad v_3=(19,\,19\sqrt3).
  \label{eq:triangle}
\end{equation}
按逆时针方向令 \(e_i=(v_{i+1}-v_i)/L\)，即三条边的单位方向向量，其中采用循环下标 \(v_4=v_1\)。取三个检测点
\begin{equation}
  s_i=v_i-1000\,e_i,
\end{equation}
并令示向度读数 \(\theta_i=\arg(e_i)+\alpha\)，误差界仍为 \(\pm1^\circ\)。

如此构造，每个角域的下边界方向 \(\theta_i-\alpha=\arg(e_i)\) 沿三角形的一条边，叉积约束 \([u(\theta_i-\alpha),\,x-s_i]\ge0\) 保留该边内侧；上边界方向 \(\theta_i+\alpha=\arg(e_i)+2\alpha\) 相对下边界向内旋转 \(2^\circ\)。第三顶点相对检测点的偏角为
\begin{equation}
  \arctan\frac{19\sqrt3}{1019}<2\alpha=\frac{\pi}{90},
\end{equation}
因此上边界约束对整个 \(T\) 是冗余的，而三条下边界半平面的交集是 \(T\)。这个三角形可以由本题的三个测向角域相交得到。

取三角形重心为真实源位置，则三次示向度均符合 \(\pm1^\circ\) 误差；源到各检测点的距离约 \(1019\) 米，落在 \(1000\)--\(1500\) 米的有效接收半径范围内，源也位于 \(1800\) 米圆域内。

等边三角形的最小包围圆圆心为其重心，半径为 \(L/\sqrt3\)：三个顶点到任意点 \(c\) 的距离平方之和等于 \(3\lVert c-g\rVert^2+L^2\)（\(g\) 为重心），其平均至少为 \(L^2/3\)，故最大距离至少为 \(L/\sqrt3\)，且在重心处恰好取到。于是
\begin{equation}
  \operatorname{diam}(T)=38\ \text{米},\qquad \rho(T)=\frac{38}{\sqrt3}\approx21.94\ \text{米}.
  \label{eq:counterexample}
\end{equation}
\(\rho(T)>\operatorname{diam}(T)/2=19\) 米，故直径 \(38\) 米的圆不能覆盖 \(T\)；事实上直径 \(40\) 米的圆（半径 \(20\) 米）同样无法覆盖。这说明"直径不超过 \(40\) 米"不能作为保证清除的充分条件，问题一的正解必须基于最小包围圆半径 \(\rho(P)\)。

最小包围圆由其直径的两端点或三个非共线边界点确定 \cite{Welzl1991}。本文的实现先检查最远点对的直径圆，若不能覆盖全部顶点，再枚举三点外接圆并保留包含全部顶点的最小者，最坏复杂度 \(O(h^4)\)。

平面中还有 \(\operatorname{diam}(P)/2\le\rho(P)\le\operatorname{diam}(P)/\sqrt3\)。若圆由两点支撑则 \(\rho(P)=\operatorname{diam}(P)/2\)；若由三点最小支撑，则圆心位于支撑三角形内。记该三角形的最大角为 \(\gamma\in[\pi/3,\pi/2]\)、最大边长为 \(\ell\le\operatorname{diam}(P)\)，则 \(\rho(P)=\ell/(2\sin\gamma)\le\operatorname{diam}(P)/\sqrt3\)。

\subsection{保证清除判据}

光学精确定位与清除只受距离约束。因此，若 \(P\) 确实包含真实目标，则从指定位置 \(q\) 保证清除的充要条件是
\begin{equation}
  d_{\max}(q;P)=\max_{x\in P}\lVert x-q\rVert=\max_i\lVert v_i-q\rVert\le r_{\mathrm{opt}} .
  \label{eq:clear-condition}
\end{equation}

对凸多边形，范数的最大值在顶点取得，故只需检查顶点。

存在至少一个保证清除位置的充要条件是 \(\rho(P)\le r_{\mathrm{opt}}\)；在本问中，将所有保证清除位置组成的集合记为 \(\mathcal C(P)\)，并用 \(\overline B(v,r)\) 表示以 \(v\) 为圆心、半径为 \(r\) 的闭圆盘，于是
\begin{equation}
  \mathcal C(P)=\bigcap_i\overline{B}(v_i,r_{\mathrm{opt}}).
  \label{eq:clear-set}
\end{equation}
当 \(\rho(P)\le r_{\mathrm{opt}}\) 时，最小包围圆圆心 \(c(P)\) 是一个合法选择，但不一定是离机器狗当前位置最近的选择。

\subsection{最近的保证清除位置}

前面给出了保证清除的充要条件 \(\rho(P)\le r_{\mathrm{opt}}\) 与全部保证清除位置 \(\mathcal C(P)\)。若该条件满足，设机器人当前位置为 \(p\)，本问另定义工程清除半径 \(\rho_{\mathrm{clr}}=19.5\) 米，相比物理阈值 \(r_{\mathrm{opt}}=20\) 米预留 \(0.5\) 米余量。将满足该工程半径要求的行动位置集合记为 \(\mathcal C_{\mathrm{clr}}(P)\)，则
\begin{equation}
  \mathcal C_{\mathrm{clr}}(P)=\{q:\max_{x\in P}\lVert q-x\rVert\le\rho_{\mathrm{clr}}\}
  =\bigcap_{i}\overline{B}(v_i,\rho_{\mathrm{clr}}).
  \label{eq:clear-proj}
\end{equation}
\(\mathcal C_{\mathrm{clr}}(P)\) 非空当且仅当 \(P\) 的最小包围圆半径不超过 \(\rho_{\mathrm{clr}}\)。当该集合非空时，目标是求欧氏投影
\begin{equation}
  q^*=\arg\min_{q\in \mathcal C_{\mathrm{clr}}(P)}\lVert q-p\rVert ,
\end{equation}
对应本次清除行动的耗时为 \(\lVert q^*-p\rVert/5+5\) 秒。

若 \(p\in \mathcal C_{\mathrm{clr}}(P)\)，在原地清除。否则最优点位于边界：若只有一个非重复圆盘约束在该处起作用，投影满足圆弧法向条件，是该圆上朝向 \(p\) 的径向投影；若至少两个不同圆约束起作用，该点属于两圆交点。枚举每个顶点圆的径向投影、每对圆的交点，逐一检验全部圆盘约束，取距 \(p\) 最近的可行点即可。该枚举给出全局最近点，枚举至多 \(O(h^2)\) 个候选、每个检查 \(h\) 个距离，简单实现 \(O(h^3)\)。

\subsection{求解结果}

\begin{table}[H]
  \centering
  \caption{问题一的主要计算结果}
  \label{tab:q1-results}
  \begingroup
  \renewcommand{\arraystretch}{1.4}
  \begin{tabularx}{\textwidth}{L{0.30\textwidth}ZC{0.20\textwidth}}
    \toprule[1.2pt]
    {\heifont\bfseries 项目} & {\heifont\bfseries 取值} & {\heifont\bfseries 备注} \\
    \midrule[0.5pt]
    边长为 38 米的等边三角形反例直径 & 38.0000 米 & 可由三个 \(2^\circ\) 角域真实形成 \\
    同三角形的最小包围圆半径 \(\rho(T)\) & 21.9393 米 & 圆心为重心 \\
    同直径圆能否覆盖 & 否 & 最小包围圆半径大于 19 米 \\
    直径 40 米圆能否覆盖 & 否 & 反例区域仍不可覆盖 \\
    保证清除的充要条件 & \(\rho(P)\le r_{\mathrm{opt}}\) & 等价于 \(\mathcal C(P)\ne\varnothing\) \\
    顶点枚举与 LP 支持函数最大差 & \(4.55\times10^{-13}\) 米 & 枚举顶点 vs 独立支持函数 \\
    最小包围圆半径差 & \(5.96\times10^{-10}\) 米 & 支撑圆枚举 vs 凸范数优化 \\
    \bottomrule[1.2pt]
  \end{tabularx}
  \endgroup
\end{table}


\section{问题二：第二检测点的两种优化模型}

第二检测点的选择可以转化为对观测后定位效果的预先比较。具体而言，对每个候选点，先分析可能获得哪些反馈，再求出各反馈下仍与观测相容的源位置区域，最后以区域的包围圆半径构造评价指标。这样，选点时既考虑两次测向的交会效果，也考虑强信号和无信号带来的位置约束。以下先给出两种模型共用的区域构造，再分别按反馈概率求期望、按全部可能反馈取最坏值，形成侧重平均效果与不利情形保证的两种选点方案。

\subsection{由观测信息构造定位区域}

区域构造的基本方法是将每次反馈转化为位置约束，再与已有信息取交。利用旋转对称性，取 \(s_1=(a,0)\)、\(\theta_1=\beta\)，其中 \(a\ge0\)、\(\beta\in[-\pi,\pi)\) 为给定参数。考察第二点 \(q\) 及其可能读数 \(\theta\) 时，令 \(s_2=q\)、\(\theta_2=\theta\)，并简记 \(d_1(g)=\|g-s_1\|\)、\(d_2(g;q)=\|g-q\|\)。在已有角域 \(\mathcal W_i\) 上加入正常接收的距离限制，得到扇形区域 \(\mathcal F_i=\mathcal W_i\cap\{g:r_{\mathrm{strong}}<\|g-s_i\|\le1500\}\)，其中 \(\mathcal F_2\) 随 \((q,\theta)\) 变化。

第一次观测将源位置限制在 \(A_1=\Omega\cap\mathcal F_1\)，令 \(K_1=\overline{A_1}\)，并要求 \(A_1\ne\varnothing\)。交集中的 \(\Omega\) 保留了目标圆域的限制，因此第一点的位置及示向度的相对方向均会影响后续选点。对任一 \(g\in A_1\)，还须存在同一个接收半径满足 \(\max\{1000,d_1(g)\}\le R\le1500\)。这一条件在第二次反馈中继续保留。

以 \(\mathsf H\)、\(\mathsf N\) 和 \(\theta\in[0,2\pi)\) 表示强信号、无信号和正常示向度反馈。强信号给出近距约束，正常反馈增加一个扇形约束；无信号则要求首次能够接收、第二次不能接收，即存在 \(R\) 满足 \(\max\{1000,d_1(g)\}\le R<d_2(g;q)\)。消去 \(R\) 后，三类反馈对应的相容位置集为
\begin{equation}
\begin{aligned}
A_{\mathsf H}(q)&=\{g\in A_1:d_2(g;q)\le r_{\mathrm{strong}}\},\\
A_{\mathsf N}(q)&=\{g\in A_1:d_2(g;q)>\max\{1000,d_1(g)\}\},\\
A_\theta(q)&=A_1\cap\mathcal F_2(q,\theta).
\end{aligned}
\label{eq:q2-pair-regions}
\end{equation}
正常分支中的两次源距均不超过1500，因而也存在满足两次接收的共同半径。记 \(\mathcal Z(q)\) 为使 \(A_z(q)\ne\varnothing\) 的全部反馈；对这些反馈，令 \(K_z(q)=\overline{A_z(q)}\)、\(\rho_z(q)=\rho(K_z(q))\)，另记 \(\rho_1=\rho(K_1)\)。

若某点距整个 \(K_1\) 均超过1500，则必然无信号，且定位区域保持不变。因此，以下在 \(\mathcal Q(a,\beta)=\{q\in\mathbb R^2:q\ne s_1,\ \operatorname{dist}(q,K_1)\le1500\}\) 内比较候选点；该约束允许检测点位于 \(\Omega\) 外。

\subsection{最小化期望定位损失的模型}

期望模型需要回答两个相互衔接的问题：首次观测后，哪些源位置更可能；在这些位置与接收半径的条件分布下，第二点的各种反馈各有多大可能。前者由贝叶斯更新确定，后者用于对定位损失加权。

\subsubsection{首次观测后的概率更新}

用 \(G\) 表示源位置的随机向量。假设观测前 \(G\) 在 \(\Omega\) 内按面积均匀分布，\(R\sim\operatorname{Unif}[1000,1500]\)，不同检测点的角误差 \(E_1,E_2\sim\operatorname{Unif}[-\alpha,\alpha]\)，且上述随机量相互独立。首次观测信息记为 \(\mathcal I_1=(s_1,\beta)\)。

对于符合方向约束的位置，均匀角误差给出的似然密度相同，但正常接收还要求 \(R\ge d_1(g)\)。因此，需要先计算距离为 \(d\) 时的接收概率
\begin{equation}
 w(d)=\Pr(R\ge d)=
 \begin{cases}
 1,&d\le1000,\\
 (1500-d)/500,&1000<d<1500,\\
 0,&d\ge1500.
 \end{cases}
\label{eq:q2-pair-weight}
\end{equation}
由贝叶斯公式，将位置先验与上述观测似然相乘，再归一化，得到首次位置后验
\begin{equation}
 p_{1,G}(g)=\frac{\mathbf1_{A_1}(g)w(d_1(g))}{C_1},
 \qquad C_1=\int_{A_1}w(d_1(g))\,\mathrm dg>0.
\label{eq:q2-pair-posterior}
\end{equation}
这里 \(C_1\) 保证密度积分为1。该更新的直观含义是，较远的位置需要较大的接收半径才能解释首次接收，因而在更新后获得较低权重。进一步，对于正后验密度的位置，有 \(R\mid G=g,\mathcal I_1\sim\operatorname{Unif}[\max\{1000,d_1(g)\},1500]\)；位置与半径在后验中由此产生关联，第二次预测须保留这一关系。

\subsubsection{第二次反馈的概率预测}

预测第二次反馈时，将相应的反馈条件加入首次观测条件，并对未知半径积分。记 \(d_\vee(g;q)=\max\{d_1(g),d_2(g;q)\}\)。两次正常接收要求 \(R\ge d_\vee\)，对应权重为 \(w(d_\vee)\)；先接收而后无信号的权重则为 \(w(d_1)-w(d_\vee)\)。结合方向信息，三类反馈的未归一化位置权重为
\begin{equation}
\begin{aligned}
 b_{\mathsf H}(g;q)&=\mathbf1_{A_{\mathsf H}(q)}(g)w(d_1(g)),\\
 b_{\mathsf N}(g;q)&=\mathbf1_{A_1}(g)\bigl[w(d_1(g))-w(d_\vee(g;q))\bigr],\\
 b_\theta(g;q)&=\frac{\mathbf1_{A_\theta(q)}(g)}{2\alpha}w(d_\vee(g;q)).
\end{aligned}
\label{eq:q2-pair-feedback-weights}
\end{equation}
其中，正常分支的 \(1/(2\alpha)\) 来自第二次测角误差的均匀密度。令 \(M_z(q)=\int_\Omega b_z(g;q)\,\mathrm dg\)，对全部位置汇总并除以首次观测的归一化常数，即得
\begin{equation}
 P_{\mathsf H}(q)=\frac{M_{\mathsf H}(q)}{C_1},\qquad
 P_{\mathsf N}(q)=\frac{M_{\mathsf N}(q)}{C_1},\qquad
 f_\Theta(\theta;q)=\frac{M_\theta(q)}{C_1}.
\label{eq:q2-pair-prediction}
\end{equation}
前两项为离散反馈概率，最后一项为正常角度的预测密度分量，其积分等于正常反馈发生的概率，三者合计为1。实际得到反馈后，将对应权重归一化，还可得到 \(p_{2,G}(g\mid z,q)=b_z(g;q)/M_z(q)\)，其中 \(M_z(q)>0\)。至此，每种反馈的发生可能性与其定位区域均已确定，可以构造选点指标。

\subsubsection{损失函数与选点目标}

本方案兼顾正常反馈的定位精度与无信号的代价。正常反馈以区域半径计损失；强信号能够在当前点完成精确定位，取零损失；无信号采用首次区域半径加附加惩罚作为失败代价，即
\begin{equation}
 L_{\mathrm E,z}(q)=
 \begin{cases}
 0,&z=\mathsf H,\\
 \rho_\theta(q),&z=\theta,\\
 \rho_1+\lambda,&z=\mathsf N,
 \end{cases}
 \qquad \lambda\ge0.
\label{eq:q2-pair-expected-loss}
\end{equation}
基础方案取 \(\lambda=0\)，使无信号的代价不小于任何正常反馈的区域半径；增大 \(\lambda\) 则提高对无信号的惩罚。这里无信号采用的是决策代价，其实际相容区域仍由式~\eqref{eq:q2-pair-regions} 给出。

按式~\eqref{eq:q2-pair-prediction} 的概率权重求平均，得到期望损失
\begin{equation}
 J_{\mathrm E}(q)
 =P_{\mathsf N}(q)(\rho_1+\lambda)
 +\int_0^{2\pi}f_\Theta(\theta;q)\rho_\theta(q)\,\mathrm d\theta.
\label{eq:q2-pair-expected-objective}
\end{equation}
第一项衡量无信号的预期代价，第二项衡量正常反馈的平均定位损失，强信号分支以零损失参与加权。给定 \((a,\beta)\) 与 \(\lambda\)，求 \(\inf_{q\in\mathcal Q(a,\beta)}J_{\mathrm E}(q)\)，最小值可达时取 \(q_{\mathrm E}^*\in\arg\min_{q\in\mathcal Q(a,\beta)}J_{\mathrm E}(q)\)。数值上，对候选点计算反馈权重和区域半径，经角度积分得到评分，再通过空间搜索与局部加密寻找近优点。

\subsection{最小化最坏剩余半径的模型}

最坏模型着重控制不利反馈造成的定位误差。对每个候选点，考察所有与首次信息相容的第二次反馈，并以其中最大的剩余半径评价该点。这样，即使不知道各情形的发生概率，也能比较不同检测点所提供的定位保证。

本方案按反馈后的实际剩余区域计损失，并对所有反馈采用同一原地定位判据。当整个区域都在当前点的光学范围内时，可精确定位，剩余损失为零；否则用区域半径衡量尚未消除的位置不确定性。因此定义
\begin{equation}
 L_{\mathrm W,z}(q)=
 \begin{cases}
 0,&d_{\max}(q;K_z(q))\le r_{\mathrm{opt}},\\
 \rho_z(q),&d_{\max}(q;K_z(q))>r_{\mathrm{opt}}.
 \end{cases}
\label{eq:q2-pair-worst-loss}
\end{equation}
由此，强信号分支的损失为零，无信号分支的损失为 \(\rho_{\mathsf N}(q)\)，正常分支则依据上述判据确定。对全部可实现反馈取上确界，再在候选点中最小化，得到
\begin{equation}
 J_{\mathrm W}(q)=\sup_{z\in\mathcal Z(q)}L_{\mathrm W,z}(q),
 \qquad J_{\mathrm W}^*=\inf_{q\in\mathcal Q(a,\beta)}J_{\mathrm W}(q).
\label{eq:q2-pair-worst-objective}
\end{equation}
最小值可达时，取 \(q_{\mathrm W}^*\in\arg\min_{q\in\mathcal Q(a,\beta)}J_{\mathrm W}(q)\)。内层取最坏值时保留所有可实现反馈，包括边界状态及退化交集；外层据此寻找最坏损失较小的检测点。数值求解采用两层搜索，内层细分正常角度区间并与强信号、无信号分支比较，外层在检测点空间逐步加密，同时检查原地定位条件。

为给出具有较好定位效果的候选区域，两种模型均可在最优值附近保留一定余量。对 \(J\in\{J_{\mathrm E},J_{\mathrm W}\}\)，记 \(J^*=\inf_{q\in\mathcal Q}J(q)\)，给定 \(\varepsilon>0\)，取 \(\mathcal Q_{J,\varepsilon}=\{q\in\mathcal Q:J(q)\le J^*+\varepsilon\}\)。其中 \(\varepsilon\) 以米计，表示相对于相应最优指标允许增加的损失。




\section{模型检验}

% 本节原稿尚未填写，合并时保留待补充。

\section{模型评价与推广}

\subsection{模型的优点}

% 本节原稿尚未填写，合并时保留待补充。

\subsection{模型的局限性}

% 本节原稿尚未填写，合并时保留待补充。

\subsection{模型的推广}

% 本节原稿尚未填写，合并时保留待补充。

\begin{thebibliography}{99}
  \bibitem{Tokekar2013} Tokekar P, Isler V. Sensor placement and selection for bearing sensors with bounded uncertainty[C]. IEEE International Conference on Robotics and Automation (ICRA), 2013. DOI: 10.1109/ICRA.2013.6630920.
  \bibitem{Gholami2015} Gholami M R, Gezici S, Wymeersch H, et al. Characterizing the worst-case position error in bearing-only target localization[C]. Workshop on Positioning, Navigation and Communication (WPNC), 2015.
  \bibitem{Welzl1991} Welzl E. Smallest enclosing disks (balls and ellipsoids)[C]. New Results and New Trends in Computer Science, LNCS 555, 1991: 359--370. DOI: 10.1007/BFb0038202.
\end{thebibliography}

\end{document}
